% Talk-to-your-Company: presentation deck (Beamer).
% Build:  latexmk -pdf main.tex   (two passes; the counter and progress bar need the .aux)
% Theme:  Triad with the "nvglobant" palette defined in beamerthemeTriad.sty beside this file.
\documentclass[aspectratio=169,11pt,t]{beamer}

\usepackage[T1]{fontenc}
\usepackage[utf8]{inputenc}
\IfFileExists{lmodern.sty}
  {\usepackage{lmodern}\usepackage{microtype}}
  {\usepackage{anyfontsize}\usepackage[expansion=false]{microtype}}

\usepackage{booktabs}
\usepackage{graphicx}
\usepackage{tikz}
\usepackage{amsmath,amssymb}

\usetheme[palette=nvglobant,dark=true]{Triad}
\usepackage{triadgraphics}

% On this palette the fills are bright greens, so type on them is dark.
\tikzset{
  tnode/.append style={text=triadPaper},
  % colour-only variants, so they can follow a size variant without resetting it
  tsoft/.style={fill=triadWash, text=triadInk, draw=triadMuted!45, line width=0.5pt},
  talt/.style={fill=triadSecondary, text=white},
  tacc/.style={fill=triadAccent, text=triadPaper},
}
\setbeamercolor{block title alerted}{fg=triadPaper,bg=triadAccent}
\setbeamercolor{block title example}{fg=white,bg=triadSecondary}
\graphicspath{{../assets/}}

% A screenshot with a hairline frame so its dark edge separates from the dark slide.
\newcommand{\shot}[2][\linewidth]{%
  {\setlength{\fboxsep}{0pt}\setlength{\fboxrule}{0.4pt}%
   \fcolorbox{triadMuted!55}{triadPaper}{\includegraphics[width=\dimexpr#1-0.8pt\relax]{#2}}}}
\setbeamerfont{title}{size=\Huge,series=\bfseries}

\title[Talk-to-your-Company]{Talk-to-your-Company}
\subtitle{A live 3D twin of your plants that you can talk to, built for one NVIDIA Jetson Thor}
\author[Yash Kavaiya]{Yash Kavaiya}
\institute[Globant Physical AI Hackathon]{Globant Physical AI Hackathon}
\date{October 2026}

\begin{document}

\begin{frame}[plain,noframenumbering]
  \titlepage
\end{frame}

% ============================================================================
\section{The problem}

\begin{frame}{Safety events go unseen because nobody can watch every camera}
  \triadlead{A plant manager with several buildings and floors is always looking at the wrong screen.}
  \begin{itemize}
    \item Near-misses, people in restricted zones and missing helmets go unreported
    \item A simple question, ``what is happening on floor 3?'', means walking there or calling someone
    \item Camera walls show video; they do not answer questions
    \item Sending factory video to a cloud is slow, costly and often not allowed
  \end{itemize}
  \triadtakeaway{The manager needs one view of the whole site, and a way to ask it things.}
\end{frame}

\begin{frame}[plain,noframenumbering]
  \triadstatement[Ask by voice. The twin flies to the right floor, shows what the cameras see, and answers out loud.]
    {A live 3D twin of your company that you can talk to.}
\end{frame}

% ============================================================================
\section{What it does}

\begin{frame}{Both plants and all six floors are live in one view}
  \begin{columns}[T,onlytextwidth]
    \begin{column}{0.60\textwidth}
      \shot{twin.jpg}
    \end{column}
    \begin{column}{0.36\textwidth}
      \triadlead{Every person and vehicle the cameras see is a marker on the right floor.}
      \begin{itemize}
        \item 2 plants, 3 floors each, 6 cameras
        \item Positions update 5 times a second
        \item Pulsing red markers for safety events
      \end{itemize}
    \end{column}
  \end{columns}
\end{frame}

\begin{frame}{The whole site is generated from one configuration file}
  \begin{columns}[T,onlytextwidth]
    \begin{column}{0.36\textwidth}
      \triadlead{Select a plant and its floors spread apart to show what is inside.}
      \begin{itemize}
        \item Plants, floors, cameras and zones
        \item Conveyors, racks, machines, pallets, offices
        \item Change the file and the twin, the rules and the assistant all follow
      \end{itemize}
    \end{column}
    \begin{column}{0.60\textwidth}
      \shot{plants.jpg}
    \end{column}
  \end{columns}
\end{frame}

\begin{frame}{You ask in plain language, by voice or by text}
  \begin{columns}[T,onlytextwidth]
    \begin{column}{0.60\textwidth}
      \shot{ask.jpg}
    \end{column}
    \begin{column}{0.36\textwidth}
      \triadlead{The answer is spoken, shown as text, and drives the twin.}
      \begin{itemize}
        \item ``Give me a status of both plants''
        \item ``Show me Plant 2, floor 3'' moves the camera there
        \item Out-of-scope questions get an honest ``I can't see that''
      \end{itemize}
    \end{column}
  \end{columns}
\end{frame}

\begin{frame}{Four rules turn tracks into safety events on the right floor}
  \begin{columns}[T,onlytextwidth]
    \begin{column}{0.50\textwidth}
      \small
      \begin{tabular}{@{}p{24mm}p{43mm}@{}}
        \toprule
        \textbf{Event} & \textbf{Raised when} \\
        \midrule
        Near miss & a person is within 2 m of a closing forklift \\
        Restricted zone & a person stays inside for over 2 s \\
        No helmet & the vision check on a person says no \\
        Crowding & more than 4 people are within 3 m \\
        \bottomrule
      \end{tabular}
      \vskip1.2ex
      Events are stored and can be asked about by plant, floor, type and time.
    \end{column}
    \begin{column}{0.46\textwidth}
      \shot{safety.jpg}
    \end{column}
  \end{columns}
\end{frame}

\begin{frame}{Every floor has a live camera view with detections boxed}
  \begin{columns}[T,onlytextwidth]
    \begin{column}{0.52\textwidth}
      \shot[0.96\linewidth]{camera.jpg}
    \end{column}
    \begin{column}{0.44\textwidth}
      \triadlead{Focus a floor and its camera appears beside the twin.}
      \begin{itemize}
        \item The same detection that places the marker draws the box
        \item One stream is shown at a time, so it stays light
        \item Nothing is published beyond the one app port
      \end{itemize}
    \end{column}
  \end{columns}
\end{frame}

\begin{frame}{Any camera can be fed from a recording or a webcam}
  \begin{columns}[T,onlytextwidth]
    \begin{column}{0.36\textwidth}
      \triadlead{For a live demo, drop your own footage into a floor.}
      \begin{itemize}
        \item Recorded mp4 files, looped as live streams
        \item A webcam shared from the browser, or attached to the machine
        \item ``Share recording'': pick a video file and it is analysed on the spot
      \end{itemize}
    \end{column}
    \begin{column}{0.60\textwidth}
      \shot{share.jpg}
    \end{column}
  \end{columns}
\end{frame}

\begin{frame}{The alert names the person, and one question gives the full picture}
  \begin{columns}[T,onlytextwidth]
    \begin{column}{0.40\textwidth}
      \shot[0.84\linewidth]{who.jpg}
    \end{column}
    \begin{column}{0.56\textwidth}
      \triadlead{``Who is this person?''}
      \begin{itemize}
        \item Name, role, shift, supervisor and training
        \item When they entered, when the alert fired, how long they stayed
        \item Whether they were authorised for that zone
        \item Where they are now, and their other incidents today
      \end{itemize}
      \vskip0.6ex
      {\small\color{triadMuted}Employees are fictional. Identity comes from a simulated badge feed, not from face recognition.\par}
    \end{column}
  \end{columns}
\end{frame}

\begin{frame}{One more sentence writes the incident report}
  \begin{columns}[T,onlytextwidth]
    \begin{column}{0.56\textwidth}
      \triadlead{``Write the incident report.''}
      \begin{itemize}
        \item What, where, when, who, severity and recommended action
        \item Filled from the stored event, so it cannot contain invented facts
        \item The camera snapshot from the moment of the alert is attached
      \end{itemize}
      \triadtakeaway{From alert to written report without leaving the twin.}
    \end{column}
    \begin{column}{0.40\textwidth}
      \hfill\shot[0.84\linewidth]{report.jpg}
    \end{column}
  \end{columns}
\end{frame}

% ============================================================================
\section{How it works}

\begin{frame}{One process on one device takes video to a spoken answer}
  \centering
  \begin{tikzpicture}[x=0.97mm,y=1mm]
    % perception lane
    \node[tlabel, anchor=west] at (-9,8) {PERCEPTION, 6 STREAMS};
    \node[tnode sm, tsoft] (cam) at (0,0) {Cameras};
    \node[tnode sm] (ing) at (24,0) {Ingest};
    \node[tnode sm] (det) at (48,0) {Detector};
    \node[tnode sm] (trk) at (72,0) {Track + map};
    \node[tnode sm] (rul) at (96,0) {Rules};
    \node[tnode sm, talt] (sto) at (120,0) {Event store};
    % voice lane
    \node[tlabel, anchor=west] at (-9,-14) {VOICE};
    \node[tnode sm, tsoft] (mic) at (0,-22) {Mic};
    \node[tnode sm] (asr) at (24,-22) {Speech in};
    \node[tnode sm, tacc] (agt) at (48,-22) {Agent + tools};
    \node[tnode sm] (tts) at (72,-22) {Speech out};
    \node[tnode sm, tsoft] (web) at (120,-22) {Browser twin};
    \draw[tarrow] (cam)--(ing); \draw[tarrow] (ing)--(det); \draw[tarrow] (det)--(trk);
    \draw[tarrow] (trk)--(rul); \draw[tarrow] (rul)--(sto);
    \draw[tarrow] (mic)--(asr); \draw[tarrow] (asr)--(agt); \draw[tarrow] (agt)--(tts);
    \draw[tarrow] (tts)--(web) node[tlabel, midway, below]{answer and speech};
    \draw[tarrow] (sto)--(web) node[tlabel, midway, right]{state, events};
    \draw[tarrow soft] (sto.south west) to[out=215,in=20] (agt.north east);
    \node[tlabel] at (64,-7) {tools read events and live state};
  \end{tikzpicture}
  \vskip2.4ex
  \raggedright
  \begin{itemize}
    \item One vision-language model serves the conversation and the ``look at this frame'' tool
    \item Seven tools: events, live state, look, focus, show event, identify person, write report
  \end{itemize}
\end{frame}

\begin{frame}{The planned model stack leaves headroom inside 15 GB}
  \begin{columns}[T,onlytextwidth]
    \begin{column}{0.60\textwidth}
      \begin{tikzpicture}
        \begin{axis}[triad hbars, width=58mm, height=0.58\textheight,
          symbolic y coords={Headroom,App and buffers,Speech synthesis,Detector,Speech recognition,Vision-language model},
          ytick=data, xlabel={Memory budget (GB)}, xmin=0, xmax=8]
          \addplot[triad fill a] coordinates
            {(2,Headroom)(2,App and buffers)(0.3,Speech synthesis)(1,Detector)(1.5,Speech recognition)(6,Vision-language model)};
        \end{axis}
      \end{tikzpicture}
    \end{column}
    \begin{column}{0.36\textwidth}
      \triadlead{12.8 of 15 GB budgeted.}
      \begin{itemize}
        \item 3B-class vision-language model, half precision
        \item Whisper for speech in, Piper for speech out
        \item YOLO exported to TensorRT, one batch for all streams
      \end{itemize}
      \vskip0.6ex
      {\small\color{triadMuted}This is the plan. The figures still have to be measured on the Jetson.\par}
    \end{column}
  \end{columns}
\end{frame}

\begin{frame}{Every heavy component has a mock, so the demo cannot die}
  \centering
  \begin{tabular}{@{}lll@{}}
    \toprule
    \textbf{Component} & \textbf{Real} & \textbf{Mock} \\
    \midrule
    Perception & video, detector, tracker & simulated people, scripted scenarios \\
    Assistant & vision-language model & keyword router calling the same tools \\
    Speech in & Whisper & fixed question \\
    Speech out & Piper & chime \\
    \bottomrule
  \end{tabular}
  \vskip2ex
  \raggedright
  \begin{itemize}
    \item A backend that fails to load falls back to its mock at startup, and the app says which is running
    \item The whole app runs on a laptop without a GPU, which is how it was built and tested
  \end{itemize}
  \triadtakeaway{If a model fails on stage, the twin and the conversation keep working.}
\end{frame}

% ============================================================================
\section{Status and demo}

\begin{frame}{The software is proven on a laptop; the Jetson run is the next step}
  \triadmetricrow
    {\triadmetric{77}{automated tests passing}}
    {\triadmetric[triadPrimary]{6}{camera streams detected together on a laptop CPU}}
    {\triadmetric[triadMuted]{0}{runs on a Jetson so far}}
  \vskip0.9em
  \begin{columns}[T,onlytextwidth]
    \begin{column}{0.48\textwidth}
      \textbf{Shown working}
      \begin{itemize}
        \item Twin, events, conversation, report
        \item Real detection, speech in and speech out on CPU
        \item Webcam and shared recording
      \end{itemize}
    \end{column}
    \begin{column}{0.48\textwidth}
      \textbf{Still to prove on the Jetson}
      \begin{itemize}
        \item Hardware video decode and TensorRT
        \item 3B model picks tools, 13 of 15
        \item 10 fps on 6 streams, voice under 2.5 s
      \end{itemize}
    \end{column}
  \end{columns}
\end{frame}

\begin{frame}{The demo is seven steps, from overview to incident report}
  \begin{columns}[T,onlytextwidth]
    \begin{column}{0.48\textwidth}
      \begin{enumerate}
        \item Twin loads: six streams live, metrics visible
        \item ``Give me a status of both plants''
        \item ``Show me Plant 2, floor 3''
        \item ``Any safety issues in the last ten minutes?''
      \end{enumerate}
    \end{column}
    \begin{column}{0.48\textwidth}
      \begin{enumerate}
        \setcounter{enumi}{4}
        \item Click the restricted-zone alert on Plant 1, floor 2
        \item ``Who is this person?''
        \item ``Write the incident report''
      \end{enumerate}
    \end{column}
  \end{columns}
  \triadtakeaway{Then one unscripted question from a judge.}
\end{frame}

\begin{frame}{Why it matters: answers at the speed of a question, with data that stays on site}
  \triadlead{The value is the time between something happening and someone knowing.}
  \begin{itemize}
    \item \textbf{Safety:} events are seen, named and written up in the moment, not at the end of a shift
    \item \textbf{Operations:} one person can ask about any floor without walking to it
    \item \textbf{Privacy and cost:} audio and video are processed on one device, with no cloud calls at runtime
    \item \textbf{Fit:} a new site is a new configuration file, not a new project
  \end{itemize}
  \triadtakeaway{Built to run on a single NVIDIA Jetson Thor, so it can sit inside the plant.}
\end{frame}

\begin{frame}[plain,noframenumbering]
  \triadstatement[Talk-to-your-Company. Globant Physical AI Hackathon, on NVIDIA Jetson Thor.]
    {Ask your company a question. It shows you, and it tells you.}
\end{frame}

\end{document}
